\documentclass[12pt]{article}
\usepackage{enumitem}
\usepackage{geometry}
\usepackage{amsmath, amssymb}
\geometry{margin=1in}
\begin{document}
\begin{center}
Astronomy Assessment 1 \
Total Marks: 40 \
Answer all questions.
\end{center}

\section*{Multiple Choice (20 marks)}
\begin{enumerate}[label=\arabic*.]
\item What drives differentiation?
\begin{enumerate}[label=(\alph*)]
    \item Spontaneous emission from radioactive atoms.
    \item The minimization of gravitational potential energy.
    \item Thermally induced collisions.
    \item Plate tectonics.
\end{enumerate}
\item Why do we see essentially the same face of the Moon at all times?
\begin{enumerate}[label=(\alph*)]
    \item because the Moon has a nearly circular orbit around the Earth
    \item because the Moon does not rotate
    \item because the other face points toward us only at new moon when we can't see the Moon
    \item because the Moon's rotational and orbital periods are equal
\end{enumerate}
\item Previous IAAC rounds featured Proxima/Alpha Centauri as closes star(system) to the Earth. Which one is the second closest star(system)?
\begin{enumerate}[label=(\alph*)]
    \item Wolf 359
    \item Sirius
    \item 61 Cygni
    \item Barnard's Star
\end{enumerate}
\item Why is the sky blue?
\begin{enumerate}[label=(\alph*)]
    \item Because the molecules that compose the Earth's atmosphere have a blue-ish color.
    \item Because the sky reflects the color of the Earth's oceans.
    \item Because the atmosphere preferentially scatters short wavelengths.
    \item Because the Earth's atmosphere preferentially absorbs all other colors.
\end{enumerate}
\item Which of the following characteristics would not necessarily suggest that a rock we found is a meteorite.
\begin{enumerate}[label=(\alph*)]
    \item It has a fusion crust
    \item It contains solidified spherical droplets
    \item It is highly processed
    \item It has different elemental composition than earth
\end{enumerate}
\item The sky is blue because
\begin{enumerate}[label=(\alph*)]
    \item the Sun mainly emits blue light.
    \item the atmosphere absorbs mostly blue light.
    \item molecules scatter red light more effectively than blue light.
    \item molecules scatter blue light more effectively than red light.
\end{enumerate}
\item Approximately how old is the surface of Venus?
\begin{enumerate}[label=(\alph*)]
    \item 750 million years.
    \item 2 billion years.
    \item 3 billion years.
    \item 4.5 billion years.
\end{enumerate}
\item Why doesn't Venus have seasons like Mars and Earth do?
\begin{enumerate}[label=(\alph*)]
    \item Its rotation axis is nearly perpendicular to the plane of the Solar System.
    \item It does not have an ozone layer.
    \item It does not rotate fast enough.
    \item It is too close to the Sun.
\end{enumerate}
\item What is the so-called bolometric luminosity in astronomy?
\begin{enumerate}[label=(\alph*)]
    \item The luminosity integrated over vertically polarized wavelengths.
    \item The luminosity integrated over horizontally wavelengths.
    \item The luminosity integrated over visible wavelengths.
    \item The luminosity integrated over all wavelengths.
\end{enumerate}
\item Which is not a similarity between Saturn and Jupiter's atmospheres?
\begin{enumerate}[label=(\alph*)]
    \item a composition dominated by hydrogen and helium
    \item the presence of belts zones and storms
    \item an equatorial wind speed of more than 900 miles per hour
    \item significant "shear" between bands of circulation at different latitudes
\end{enumerate}
\end{enumerate}

\section*{True / False (20 marks)}
\textit{Answer True or False and justify briefly.}
\begin{enumerate}[label=\arabic*.]
\item True or False: The correct answer to 'From shortest to longest wavelength which of the following correctly orders the different categories of electromagnetic radiation?' is 'gamma rays X rays ultraviolet visible light infrared radio'.
\item True or False: The correct answer to 'Which is not an essential requirement for life as we know it?' is 'A source of organic molecules'.
\item True or False: The correct answer to 'Why do we look for water-ice in craters at Mercury's pole?' is 'These craters contain the only permanently shadowed regions on Mercury'.
\item True or False: The correct answer to 'Which of these has NOT been one of the main hypotheses considered for the origin of the Moon?' is 'The Moon split from the Earth due to tidal forces.'.
\item True or False: The correct answer to 'The distance between the Earth and the star Altair is one million times greater than the distance between the Earth and the Sun. How far...' is '1.5 x 1014 meters'.
\item True or False: The correct answer to 'When was the telescope invented by Galileo?' is '1609'.
\item True or False: The correct answer to 'What does the astronomical term ecliptic describe?' is 'The path of the Sun in the sky throughout a year.'.
\item True or False: The correct answer to 'Sunspots are black regions that temporarily appear on the Sun. Their number highly increases every ... years. This period is also called...' is '11'.
\item True or False: The correct answer to 'We were first able to accurately measure the diameter of Pluto from:' is 'brightness measurements made during mutual eclipses of Pluto and Charon'.
\item True or False: The correct answer to 'What is true for a type-Ia ("type one-a") supernova?' is 'This type produces gamma-ray bursts.'.
\end{enumerate}

\vfill
\noindent\textit{End of Paper}
\end{document}